\scanpage{26}
under the anabelian hypothesis made, and that of dominance. (This second one is not superfluous, as one sees by taking \deleted{for $Y$ \inserted{$=G\backslash Z$, where $G$ is} an elliptic group and $Z$ a finite subgroup $\unsure{\ne 1|+1}$, and for $f,g$} \inserted{two constant maps from $X$ to $Y$!})

First of all, one checks that an elementary Artin variety does indeed embed \inserted{by an $i:Y\longrightarrow G$} in an extension of an abelian variety by a torus -- so one will have $(i\circ g)=\tau_a(i\circ f)$. \deleted{Thus \unsure{all}} \inserted{On the other hand,} \deleted{Replacing if necessary} $f$,\deleted{$g$} being dominant, there exists a \inserted{nonempty} open subset $\unsure{V|W}\subset Y$ contained in $\operatorname{Im}f$, and one will therefore have $\unsure{V|W}+a\subset Y$. Let $U\subset X$ be a nonempty open subset such that $f(U)\subset V$ \deleted{$g(U)\subset V$} \deleted{\unsure{therefore which} $\unsure{Q}$} replacing $X$ by $U$ if necessary, one is reduced to the case where $f$\deleted{,$g$} sends $X$ into an open subset $V$ of $Y$ such that $V+a\subset Y$, and $g\circ i=\tau_a(i\circ f)$. One \deleted{knows} \inserted{sees} that $\pi_1(f)$ is a composite (where $f':X\longrightarrow U$ is \deleted{\unsure{$f$}} induced by $f$ and $j:U\longrightarrow Y$ is the inclusion)
\diagram{\begin{tikzcd}[ampersand replacement=\&]\pi_1(X) \arrow[r,"\pi_1(f')"] \& \pi_1(U) \arrow[r,"\pi_1(j)"] \& \pi_1(Y)\end{tikzcd}}
whereas $\pi_1(g)$ is the composite
\diagram{\begin{tikzcd}[ampersand replacement=\&]\pi_1(X) \arrow[r,"\pi_1(f')"] \& \pi_1(U) \arrow[r,"{\pi_1(\unsure{j'|g'})}"] \& \pi_1(Y)\end{tikzcd}}
where $\deleted{\unsure{j}}$ $j':U\longrightarrow Y$ is induced by $\tau_a$.
\margin{25}
